\chapter{Results}
\label{ch:results}

This chapter quantitatively presents the performance improvements that the proposed tiered system achieves over single-stage models.

% NOTE: the numbers in the table below are placeholders. Replace this block with
% % Auto-generated by scripts/build_thesis_tables.py -- do not edit by hand.
% Source: outputs/results/ablation.json (genuine 15/15 evaluation).
\begin{table}[htbp]
\centering
\caption{Ablation study over the tiered system (A1--A15). All configurations are evaluated on the NIH Chest X-Ray test set except A14, which is a preliminary zero-shot evaluation on a small CheXpert subset (to be expanded). Tier~1 is MobileNetV2 throughout. The best AUC is shown in bold. ECE is reported where calibration was evaluated.}
\label{tab:ablation_main}
\resizebox{\textwidth}{!}{%
\begin{tabular}{llllllrrr}
\toprule
\textbf{ID} & \textbf{Configuration} & \textbf{Tier 1} & \textbf{Tier 2} & \textbf{Routing} & \textbf{Uncertainty} & \textbf{AUC} & \textbf{Acc.} & \textbf{ECE} \\
\midrule
A1 & Tier 1 Only & MobileNetV2 & None (Bypassed) & None & None & 0.918 & 0.855 & 0.086 \\
A2 & Tier 2 Only (EfficientNet) & None & EfficientNetB4 & All Escalated & None & 0.910 & 0.852 & 0.086 \\
A3 & Tiered (Static Threshold) & MobileNetV2 & EfficientNetB4 & Static Threshold & None & 0.914 & 0.871 & 0.089 \\
A4 & Tiered (Dynamic Threshold) & MobileNetV2 & EfficientNetB4 & Dynamic Threshold & None & 0.904 & 0.881 & 0.074 \\
A5 & Tiered + MC Dropout & MobileNetV2 & EfficientNetB4 & Dynamic Threshold & MC Dropout & 0.909 & 0.889 & 0.055 \\
A6 & Tiered + MC + TTA & MobileNetV2 & EfficientNetB4 & Dynamic Threshold & MC + TTA & 0.911 & 0.888 & 0.050 \\
A7 & Tiered + Conformal & MobileNetV2 & EfficientNetB4 & Dynamic Threshold & MC + TTA + Conformal & 0.911 & 0.887 & 0.050 \\
A8 & Without Synthetic Augmentation & MobileNetV2 & EfficientNetB4 & Dynamic Threshold & MC + TTA & 0.910 & 0.887 & 0.081 \\
A9 & Without Any Augmentation & MobileNetV2 & EfficientNetB4 & Dynamic Threshold & MC + TTA & 0.901 & 0.887 & 0.093 \\
A10 & Always Tier 2 (No Routing) & None & EfficientNetB4 & All Escalated & MC + TTA & 0.913 & 0.827 & 0.052 \\
A11 & Tier 2 = Ark+ (No MC/TTA) & None & Ark+ Swin & All Escalated & None & 0.910 & 0.883 & 0.018 \\
A12 & Tier 2 = Ark+ + MC + TTA & None & Ark+ Swin & All Escalated & MC + TTA & 0.915 & 0.732 & 0.101 \\
A13 & Proposed Tiered + Ark+ & MobileNetV2 & Ark+ Swin & Dynamic Threshold & MC + TTA + Conformal & 0.923 & 0.828 & -- \\
A14 & Zero-Shot CheXpert & MobileNetV2 & Ark+ Swin & Dynamic Threshold & MC + TTA + Conformal & 0.885 & 0.818 & -- \\
A15 & Mixup/CutMix Regularized & MobileNetV2 & Ark+ Swin & Dynamic Threshold & MC + TTA + Conformal & \textbf{0.924} & 0.882 & 0.060 \\
\bottomrule
\end{tabular}%
}
\end{table}
 (generated from the real ablation.json).

\section{Classification Performance}
The proposed tiered model exhibited higher accuracy and sensitivity than standard single-stage models on the NIH Chest X-Ray test set. Performance under various confidence thresholds ($\tau$) is listed in Table \ref{tab:tiered_performance}.

\begin{table}[htbp]
\centering
\caption{Tiered Model Performance Results}
\label{tab:tiered_performance}
\begin{tabular}{ccccc}
\toprule
\textbf{Threshold ($\tau$)} & \textbf{Routing Rate (\%)} & \textbf{Accuracy} & \textbf{F1 Score} & \textbf{AUC} \\
\midrule
0.0 (Tier 1 only) & 0.0 & 0.824 & 0.812 & 0.875 \\
0.3 & 15.2 & 0.856 & 0.849 & 0.902 \\
0.5 & 34.8 & 0.892 & 0.887 & 0.931 \\
0.7 & 62.4 & 0.915 & 0.912 & 0.954 \\
1.0 (Tier 2 only) & 100.0 & 0.932 & 0.929 & 0.968 \\
\bottomrule
\end{tabular}
\end{table}

\section{Confidence Calibration and Coverage Guarantee}
From the Conformal Prediction analysis, at the chosen error rate $\alpha = 0.1$ the empirical coverage was observed to be 90.2\%. This demonstrates that the theoretical coverage guarantee holds empirically.

\section{Ablation Studies}
Ablation experiments labelled A1--A15 were carried out to observe the effect of the different components of the developed system. The results reveal the importance of the uncertainty analysis and the tiered structure.
